% ===========================================================================
%  Annexe — Glossaire des notions
% ===========================================================================
\chapter{Glossaire des notions}

Les mots de l'année, dans l'ordre alphabétique. Chacun porte, entre
parenthèses, le numéro du chapitre où la notion est construite — c'est là
qu'il faut retourner si la définition ne suffit pas.

% Mise en forme locale à cette annexe
\newcommand{\glolettre}[1]{%
  \par\vspace{0.9em}\noindent
  {\sffamily\bfseries\large\color{bmtAccent}#1}%
  \par\vspace{-0.35em}\noindent
  {\color{bmtGrisClair}\rule{\linewidth}{0.8pt}}\par\vspace{0.45em}}

\newcommand{\glo}[3]{%
  \par\vspace{0.32em}\noindent
  {\sffamily\bfseries\small #1}%
  \hspace{0.35em}{\sffamily\footnotesize\color{bmtGris}(ch. #2)}%
  \hspace{0.45em}\ignorespaces #3\par}

% ---------------------------------------------------------------------------
\glolettre{A}

\glo{Antécédent}{3}{Nombre $x$ dont l'image par $f$ vaut $y$. Un nombre peut
avoir plusieurs antécédents, ou aucun.}

\glo{Arrangement}{10}{Liste \emph{ordonnée} de $p$ objets distincts choisis
parmi $n$. Leur nombre est $\arrang{p}{n} = \frac{n!}{(n-p)!}$ ; c'est le
décompte d'un tirage successif sans remise.}

\glo{Asymptote horizontale}{4}{Droite $y = \ell$ dont la courbe se rapproche
à l'infini, lorsque $\limpinf f(x) = \ell$.}

\glo{Asymptote oblique}{4}{Droite $y = ax+b$ telle que
$\lim\left[f(x)-(ax+b)\right] = 0$.}

\glo{Asymptote verticale}{4}{Droite $x = a$ dont la courbe se rapproche
lorsque $f(x)$ tend vers l'infini en $a$.}

\glo{Automatisme}{—}{Calcul qui doit sortir sans papier ni réflexion. Chaque
chapitre en propose huit.}

% ---------------------------------------------------------------------------
\glolettre{B}

\glo{Branche parabolique}{4}{Comportement d'une courbe qui part à l'infini
sans suivre aucune droite. De direction $(Oy)$ si
$\frac{f(x)}{x} \to \pm\infty$, de direction $(Ox)$ si
$\frac{f(x)}{x} \to 0$.}

% ---------------------------------------------------------------------------
\glolettre{C}

\glo{Cas favorables}{11}{Nombre d'issues qui réalisent l'événement étudié.
Numérateur de la fraction de probabilité.}

\glo{Centre de symétrie}{3}{Point $\Omega(a\,;\,b)$ tel que
$f(a-h)+f(a+h) = 2b$ pour tout $h$ admissible.}

\glo{Coefficient multiplicateur}{12}{Nombre $1+t$ par lequel on multiplie
une valeur pour lui appliquer une évolution de taux $t$.}

\glo{Combinaison}{10}{Partie de $p$ objets choisis parmi $n$, \emph{sans}
ordre. Leur nombre est $\combi{p}{n} = \frac{n!}{p!(n-p)!}$ ; c'est le
décompte d'un tirage simultané.}

\glo{Courbe représentative}{3}{Ensemble des points $M(x\,;\,f(x))$ pour $x$
dans l'ensemble de définition.}

\glo{Cramer (méthode de)}{2}{Résolution d'un système linéaire par
déterminants : $x = \frac{D_{x}}{D}$, $y = \frac{D_{y}}{D}$, et de même pour
$z$ à l'ordre trois.}

% ---------------------------------------------------------------------------
\glolettre{D}

\glo{Dénombrement}{10}{Art de compter le nombre de résultats possibles d'une
expérience, préalable indispensable au calcul d'une probabilité.}

\glo{Dérivée}{5}{Fonction $f'$ qui donne, en chaque point, le coefficient
directeur de la tangente à la courbe.}

\glo{Déterminant}{2}{Nombre associé à un système linéaire, qui décide de
l'existence et de l'unicité de la solution.}

\glo{Discriminant}{1}{Nombre $\Delta = b^{2}-4ac$ qui distingue les trois cas
d'une équation du second degré.}

\glo{Droite d'ajustement}{12}{Droite qui résume un nuage de points et permet
d'estimer une valeur non observée.}

% ---------------------------------------------------------------------------
\glolettre{E}

\glo{Ensemble de définition}{3}{Ensemble des réels pour lesquels
l'expression de $f$ a un sens. On le note $\Df$.}

\glo{Équiprobabilité}{11}{Situation où toutes les issues ont la même chance.
Signalée par « indiscernables au toucher » ou « tirages équiprobables ».}

\glo{Événement}{11}{Partie de l'univers. Il est \emph{certain} s'il est
l'univers entier, \emph{impossible} s'il est vide.}

\glo{Événement contraire}{11}{Événement $\evtc{A}$ réalisé exactement quand
$A$ ne l'est pas ; $\proba{\evtc{A}} = 1-\proba{A}$.}

\glo{Exponentielle}{8}{Fonction réciproque du logarithme népérien,
strictement positive et strictement croissante sur $\R$.}

\glo{Extremum}{5}{Maximum ou minimum. C'est une \emph{ordonnée}, atteinte en
une abscisse — les deux ne se confondent pas.}

\glo{Extrapolation}{12}{Estimation faite au-delà des valeurs observées. Elle
suppose que la tendance se poursuit ; ce n'est jamais une certitude.}

% ---------------------------------------------------------------------------
\glolettre{F}

\glo{Factorielle}{10}{$n! = n\times(n-1)\times\cdots\times1$, avec la
convention $0! = 1$.}

\glo{Forme indéterminée}{4}{Écriture — $\infty-\infty$, $0\times\infty$,
$\frac{\infty}{\infty}$, $\frac00$ — qui ne permet pas de conclure
directement. Il faut transformer, jamais s'arrêter là.}

\glo{Fonction paire}{3}{$f(-x) = f(x)$ : la courbe admet l'axe des ordonnées
pour axe de symétrie.}

\glo{Fonction impaire}{3}{$f(-x) = -f(x)$ : la courbe admet l'origine pour
centre de symétrie.}

% ---------------------------------------------------------------------------
\glolettre{I}

\glo{Image}{3}{Nombre $f(x)$ associé à $x$. Un nombre a \emph{une seule}
image.}

\glo{Intégrale}{9}{Nombre $\int_{a}^{b} f(x)\dx = F(b)-F(a)$, où $F$ est une
primitive de $f$.}

% ---------------------------------------------------------------------------
\glolettre{L}

\glo{Limite}{4}{Valeur dont $f(x)$ s'approche lorsque $x$ tend vers une
borne de l'ensemble de définition.}

\glo{Logarithme népérien}{7}{Fonction $\ln$, définie sur
$\intervalleo{0}{+\infty}$, de dérivée $\frac1x$, qui transforme les
produits en sommes.}

% ---------------------------------------------------------------------------
\glolettre{M}

\glo{Mayer (méthode de)}{12}{Ajustement d'un nuage par la droite joignant
les points moyens des deux moitiés de la série.}

\glo{Moyenne}{12}{Somme des valeurs divisée par leur nombre.}

% ---------------------------------------------------------------------------
\glolettre{N}

\glo{Nombre dérivé}{5}{Valeur $f'(a)$ : le coefficient directeur de la
tangente au point d'abscisse $a$.}

\glo{Nuage de points}{12}{Représentation graphique d'une série statistique
double, un point par couple observé.}

% ---------------------------------------------------------------------------
\glolettre{P}

\glo{Permutation}{10}{Rangement de $n$ objets dans un ordre. Leur nombre est
$n!$.}

\glo{Point moyen}{12}{Point $\pmoyen(\moy x\,;\,\moy y)$ du nuage. La droite
de Mayer y passe toujours.}

\glo{Position relative}{6}{Situation de la courbe par rapport à une droite,
donnée par le signe de la différence $f(x)-(ax+b)$.}

\glo{Primitive}{9}{Fonction $F$ telle que $F' = f$. Il y en a une infinité,
toutes égales à une constante près.}

\glo{Probabilité}{11}{Nombre compris entre $0$ et $1$ mesurant la chance
qu'un événement se réalise.}

% ---------------------------------------------------------------------------
\glolettre{S}

\glo{Sarrus (règle de)}{2}{Procédé de calcul d'un déterminant d'ordre trois,
et de cet ordre seulement.}

\glo{Série statistique double}{12}{Liste de couples
$\left(x_{i}\,;\,y_{i}\right)$ observés sur une même population.}

\glo{Système linéaire}{2}{Ensemble d'équations du premier degré devant être
vérifiées simultanément.}

% ---------------------------------------------------------------------------
\glolettre{T}

\glo{Tableau de variation}{5}{Tableau qui résume, pour chaque intervalle, le
signe de $f'$ et le sens de variation de $f$, avec les limites et les
extremums.}

\glo{Tangente}{5}{Droite d'équation $y = f'(a)(x-a)+f(a)$, qui épouse la
courbe au point d'abscisse $a$.}

\glo{Taux d'évolution}{12}{Quotient $\frac{V_{1}-V_{0}}{V_{0}}$, exprimé en
pourcentage.}

\glo{Tirage simultané}{10}{Tirage « d'un seul coup », où l'ordre ne compte
pas : on dénombre avec $\combi{p}{n}$.}

\glo{Tirage successif}{10}{Tirage un par un, où l'ordre compte : avec remise
on dénombre avec $n^{p}$, sans remise avec $\arrang{p}{n}$.}

\glo{Trinôme}{1}{Expression $ax^{2}+bx+c$ avec $a \neq 0$.}

% ---------------------------------------------------------------------------
\glolettre{U}

\glo{Unité d'aire}{9}{Unité dans laquelle s'exprime une intégrale. Si
l'unité graphique vaut $u$ centimètres, une unité d'aire vaut $u^{2}$
centimètres carrés.}

\glo{Univers}{11}{Ensemble $\Omega$ de tous les résultats possibles d'une
expérience aléatoire.}

% ---------------------------------------------------------------------------
\glolettre{V}

\glo{Valeur interdite}{4}{Réel qui annule un dénominateur : il est exclu de
l'ensemble de définition, et donne souvent une asymptote verticale.}

\glo{Variation (sens de)}{5}{Croissance ou décroissance d'une fonction sur un
intervalle donné.}
